\Needspace{13\baselineskip}
\section{A compact dual problem over Gauduchon currents}
\label{sec:weak-gauduchon-duality}

We now fix a Hermitian background and vary the density. For a smooth positive density, \cref{thm:gaud-volume-dual-formula} provides a minimizing Gauduchon metric; for a density that is only continuous or integrable, there is no smooth equation to produce one. The extremal problem becomes compact when positive $\ddc$-closed $(n-1,n-1)$-currents are allowed as tests, provided the denominator $D_\mu$ is extended to currents. For this we take the determinant of the matrix density of a current relative to its trace measure. The integral of its $n$th root is an infimum of continuous linear functionals, hence upper semicontinuous under weak convergence, and its maximum over normalized currents is attained. The smooth equation identifies this maximum with the supremum over smooth tests.

From here until \cref{sec:gaud-volume-nef-limit}, we assume $n\geq2$, and we fix a Hermitian form $\omega$ and a smooth positive probability volume form $\mathrm{d}V$ on $X$; integrability and convergence in measure refer to $\mathrm{d}V$. As in \cref{chap:products}, $\mathrm{MA}_\omega(u)=(\omega+\ddc u)^n$ for bounded $u\in\PSH(X,\omega)$, and $\mathrm{MA}^{\mathrm{ext}}_\omega(u)$ is the extended non-pluripolar Monge--Amp\`ere measure of \cref{thm:arbitrary-extended-ma}. Let $\mathscr G_\omega$ be the compact convex set of normalized Gauduchon currents of \cref{sec:gauduchon-linear-duality}, that is, of positive $(n-1,n-1)$-currents $T$ with $\ddc T=0$ and $\int_X\omega\wedge T=1$. Finally, let $F\in C(X,\mathbb R)$.

For $T\in\mathscr G_\omega$, the trace measure $\sigma_T=\omega\wedge T$ is a probability measure, and all coefficient measures of $T$ are absolutely continuous with respect to it. In a unitary frame for $\omega$, let $B$ be the matrix of a smooth real $(1,1)$-form $\beta$ relative to $\omega$. The coefficient densities of $T$ form a measurable positive semidefinite Hermitian matrix $P_T$, determined $\sigma_T$-almost everywhere by
\[
 \beta\wedge T=\frac1n\operatorname{tr}(B P_T)\,\mathrm{d}\sigma_T,
 \qquad \operatorname{tr}P_T=n\quad\sigma_T\text{-almost everywhere}.
\]
The matrix $P_T$ transforms by unitary conjugation when the frame changes, so its determinant is well defined. Define
\begin{equation}\label{eq:weak-gaud-functional}
 f_F=\left(\frac{\mathrm{e}^F\mathrm{d}V}{\omega^n}\right)^{1/n},\qquad
 \Gamma_F(T)=f_F(\det P_T)^{1/n}\sigma_T,\qquad
 \mathfrak D_F(T)=\Gamma_F(T)(X).
\end{equation}
These quantities are finite, since the arithmetic--geometric mean inequality gives $(\det P_T)^{1/n}\leq1$; the measure $\Gamma_F(T)$ may vanish for some currents. The same definitions apply to every positive $(n-1,n-1)$-current $T$ of finite mass, with the finite measure $\sigma_T=\omega\wedge T$ and $\operatorname{tr}P_T=n$, and $\Gamma_F$ is positively homogeneous of degree one in $T$. If the coefficient measures of $T$ are written instead as a matrix density $P'$ relative to another positive measure $\sigma'$ with respect to which they are absolutely continuous, then $(\det P_T)^{1/n}\sigma_T=(\det P')^{1/n}\sigma'$, by homogeneity of the determinant.

For a smooth positive $T=\gamma^{n-1}$, we may therefore use the reference measure $\omega^n$. The coefficient matrix of $T$ relative to $\omega^n$ is the cofactor matrix of the matrix of $\gamma$, and its determinant is $(\gamma^n/\omega^n)^{n-1}$. The definition therefore gives
\[
 \mathfrak D_F(\gamma^{n-1})
 =\int_X\left(\frac{\mathrm{e}^F\mathrm{d}V}{\gamma^n}\right)^{1/n}\gamma^n.
\]
Thus $\mathfrak D_F$ extends the denominator $D_\mu$ of \cref{def:gaud-variational-volume} to currents, with $\mu$ replaced by the possibly unnormalized density $\mathrm{e}^F\mathrm{d}V$.

\begin{lemma}\label{lem:weak-gaud-determinant-dual}
For every positive $(n-1,n-1)$-current $T$ of finite mass on $X$,
\[
 \mathfrak D_F(T)
   =\inf_{\substack{\beta>0\ \mathrm{continuous}\\\beta^n=\mathrm{e}^F\mathrm{d}V}}
                         \int_X\beta\wedge T.
\]
In particular, $\mathfrak D_F$ is concave, positively homogeneous and weakly upper semicontinuous on $\mathscr G_\omega$, and its value does not depend on the auxiliary frame or on the metric used to express the matrix density.
\end{lemma}
\begin{proof}
Let $I$ be the identity matrix of size $n$. For any positive semidefinite Hermitian matrix $P$ of that size,
\[
 (\det P)^{1/n}
      =\inf_{Q>0,\ \det Q=1}\frac1n\operatorname{tr}(QP).
\]
The inequality $\leq$ is the arithmetic--geometric mean inequality. For the reverse inequality, take, for real $\varepsilon>0$,
\[
 Q_\varepsilon=
       \det(P+\varepsilon I)^{1/n}(P+\varepsilon I)^{-1}.
\]
Its determinant is one. If $P$ is positive definite, the trace expression $\frac1n\operatorname{tr}(Q_\varepsilon P)$ converges to $(\det P)^{1/n}$ as $\varepsilon\downarrow0$. If $P$ is singular, it converges to zero: each ratio of an eigenvalue of $P$ to that eigenvalue plus $\varepsilon$ is bounded, while the prefactor tends to zero. When $\operatorname{tr}P=n$, the trace expression is at most $1+\varepsilon$.

We apply this formula to $P_T$ and multiply by $f_F$. Every continuous positive $\beta$ in the asserted infimum has matrix $f_FQ$ with $\det Q=1$, so it gives an upper bound for $\mathfrak D_F(T)$. Conversely, for fixed $\varepsilon>0$ the measurable field $Q_\varepsilon$ formed with $P=P_T$ is bounded above and below by positive constants. It can be approximated in $L^1(\sigma_T)$ by continuous positive definite Hermitian fields with bounds of the same type. One way to construct them is to apply Lusin's theorem in finitely many bundle charts, extend the continuous fields from the resulting compact sets, and glue with a partition of unity; spectral truncation to a fixed positive interval preserves the bounds. Finally we divide each approximating field by the $n$th root of its determinant. This normalization is uniformly continuous on fields with spectrum in a fixed compact positive interval, so the normalized fields still converge in $L^1(\sigma_T)$ and have determinant one. Their trace integrals converge to that of $Q_\varepsilon$.

Now let $\varepsilon\downarrow0$. The bound $1+\varepsilon$ permits dominated convergence, and we obtain equality of the infimum and the integral. For a current of arbitrary positive finite trace mass we use homogeneity, and the zero current is immediate. The infimum is taken over functionals $T\mapsto\int_X\beta\wedge T$ that are linear and weakly continuous, which proves concavity and upper semicontinuity. It also expresses $\mathfrak D_F(T)$ without a matrix frame or an auxiliary metric, which proves the final assertion.
\end{proof}

Let $B_0(F)$ be the normalization scalar of \cref{sec:normalization-scalar} for the background $\omega$ and the reference potential zero. Thus the equation
$\mathrm{MA}_\omega(u)=\mathrm{e}^{B_0(F)+F}\mathrm{d}V$
has a bounded solution $u\in\PSH(X,\omega)$, and $B_0(F)$ is the only constant with this property.

\begin{theorem}\label{thm:weak-gaud-compact-duality}\index{Gauduchon duality!weak compact formulation}
For every $F\in C(X,\mathbb R)$,
\begin{equation}\label{eq:weak-gaud-max-formula}
 M(F):=\max_{T\in\mathscr G_\omega}\mathfrak D_F(T)
                      =\mathrm{e}^{-B_0(F)/n}>0.
\end{equation}
The maximum equals the supremum over the smooth strictly positive members of $\mathscr G_\omega$. The set $\mathscr M(F)$ of maximizing currents is nonempty, compact and convex.

If $F$ is smooth, there is exactly one maximizing current, namely the normalized $(n-1)$st power of the Gauduchon metric conformal to the smooth solution form $\omega+\ddc u$ of $\mathrm{MA}_\omega(u)=\mathrm{e}^{B_0(F)+F}\mathrm{d}V$.
\end{theorem}
\begin{proof}
Upper semicontinuity of $\mathfrak D_F$ on the compact set $\mathscr G_\omega$ gives attainment, and testing a smooth positive member gives a positive value. Concavity makes the set of maximizers convex, and upper semicontinuity makes it closed.

Suppose first that $F$ is smooth, and let $\Omega=\omega+\ddc u>0$ be the smooth solution form with $\Omega^n=\mathrm{e}^{B_0(F)+F}\mathrm{d}V$. The form $\beta=\mathrm{e}^{-B_0(F)/n}\Omega$ is an admissible test in the preceding lemma, so Stokes' theorem for currents gives, for every $T\in\mathscr G_\omega$,
\[
 \mathfrak D_F(T)\leq\int_X\beta\wedge T
  =\mathrm{e}^{-B_0(F)/n}\int_X\omega\wedge T=\mathrm{e}^{-B_0(F)/n}.
\]
If $\gamma$ is the Gauduchon metric conformal to $\Omega$, equality holds in the pointwise matrix inequality for $T$ proportional to $\gamma^{n-1}$; normalizing $T$ so that $\int_X\omega\wedge T=1$, we obtain the formula and its attainment by a smooth test.

For uniqueness, let $T$ be a maximizer. Equality in the integrated inequalities forces equality in the matrix arithmetic--geometric mean inequality $\sigma_T$-almost everywhere. Since the matrix of $\beta$ is positive definite and $\operatorname{tr}P_T=n$, this means that the coefficient measures of $T$ are a nonnegative scalar measure times the smooth coefficients of $\gamma^{n-1}$. Thus $T=a\gamma^{n-1}$, where $a$ is a positive measure, regarded as a distribution. The equation $\ddc(a\gamma^{n-1})=0$ is the adjoint of a scalar elliptic equation, so elliptic regularity makes $a$ smooth. Since $\gamma$ is Gauduchon, \cref{prop:adjoint-gauduchon} identifies the kernel of this adjoint operator with the constants: the trace operator $v\mapsto n\ddc v\wedge\gamma^{n-1}/\gamma^n$ has the constants as kernel and index zero, and integration against $\gamma^n$ annihilates its range. Thus $a$ is constant, and the normalization fixes the constant. This proves uniqueness.

For continuous $F,G$, the pointwise definition gives
\[
 \mathrm{e}^{-\|F-G\|_\infty/n}\mathfrak D_G(T)
 \leq\mathfrak D_F(T)
 \leq \mathrm{e}^{\|F-G\|_\infty/n}\mathfrak D_G(T).
\]
The same bounds hold for the maxima, and for the suprema over smooth positive members. We approximate $F$ uniformly by smooth functions. Since $B_0$ is $1$-Lipschitz for the uniform norm by \cref{prop:relative-normalization-order}, the formula just proved for smooth data gives \eqref{eq:weak-gaud-max-formula} in the limit, including the equality with the smooth supremum. Note that only $F$ is approximated here: some normalized Gauduchon currents are not limits of smooth positive ones, as \cref{ex:gauduchon-current-smoothing-obstruction} shows.
\end{proof}

\section{Supporting probabilities represented by maximizing currents}
\label{sec:gaud-current-adjoints}

By \cref{thm:weak-gaud-compact-duality}, $B_0=-n\log M$, and $M(F)$ is the maximum of the functionals $T\mapsto\mathfrak D_F(T)=\int_X\mathrm{e}^{F/n}\,\mathrm{d}\Gamma_0(T)$, each of which is linear in $\mathrm{e}^{F/n}$. The supporting probabilities of $B_0$ should therefore be the normalized determinant measures of maximizing currents, and we show that they are exactly these. The point that needs proof is the continuity of the determinant measures: on the set of maximizers, upper semicontinuity of the total determinant integral becomes equality. Recall from \cref{sec:normalization-supports} that the set $\mathcal A_0(F)$ of \emph{supporting probabilities} of $B_0$ at $F\in C(X)$ consists of the Borel probability measures $\nu$ with $B_0(G)\leq B_0(F)-\int_X(G-F)\,\mathrm{d}\nu$ for all $G\in C(X)$; by \cref{thm:weak-adjoint-probabilities} it is nonempty, convex and weakly compact.

For $T\in\mathscr M(F)$ put
\[
 \nu_{T,F}=\Gamma_F(T)/M(F).
\]
This is a probability measure.

\begin{lemma}\label{lem:gaud-maximizer-measure-continuity}
Suppose $T_j\to T$ weakly in $\mathscr G_\omega$ and $\mathfrak D_F(T_j)\to M(F)$. Then $T\in\mathscr M(F)$ and $\Gamma_F(T_j)\rightharpoonup\Gamma_F(T)$. On $\mathscr M(F)$ the map $T\mapsto\Gamma_F(T)$ is affine and continuous.
\end{lemma}
\begin{proof}
Upper semicontinuity gives $\mathfrak D_F(T)\geq M(F)$, hence equality. The measures $\Gamma_F(T_j)$ have bounded mass; let $\Gamma'$ be a weak cluster point, whose mass is $M(F)$. For a continuous real function $h$, choose $\varepsilon>0$ with $1\pm\varepsilon h>0$. Since $\Gamma_{F+n\log(1\pm\varepsilon h)}=(1\pm\varepsilon h)\Gamma_F$, upper semicontinuity of $\mathfrak D$ for these two continuous logarithmic densities gives
\[
 \int_X(1\pm\varepsilon h)\,\mathrm{d}\Gamma'
       \leq\int_X(1\pm\varepsilon h)\,\mathrm{d}\Gamma_F(T).
\]
Since both measures have mass $M(F)$, the two signs give equality of the integrals of $h$. Thus every cluster point equals $\Gamma_F(T)$, which proves the convergence.

For affineness, take $T_1,T_2\in\mathscr G_\omega$ and $t\in[0,1]$, and express the coefficient measures of both currents through densities with respect to the sum of their trace measures. Concavity and homogeneity of the $n$th root of the determinant, applied to these matrix densities, give
\[
 \Gamma_F(tT_1+(1-t)T_2)
    \geq t\Gamma_F(T_1)+(1-t)\Gamma_F(T_2),\qquad 0\leq t\leq1.
\]
If $T_1$ and $T_2$ maximize, so does the current on the left, by concavity of $\mathfrak D_F$. Both measures in the inequality then have mass $M(F)$, so equality holds. This proves affineness, and continuity is the special case of the first assertion in which every $T_j$ is a maximizer.
\end{proof}

\Needspace{9\baselineskip}
\begin{theorem}\label{thm:gaud-current-all-adjoints}
For every $F\in C(X,\mathbb R)$,
\[
 \mathcal A_0(F)=\{\nu_{T,F}:T\in\mathscr M(F)\}.
\]
Consequently the determinant measures $\Gamma_F(T)$ of all maximizing currents are absolutely continuous with respect to $\mathrm{d}V$. (This concerns the measures $\Gamma_F(T)$, not the currents $T$ themselves.)
\end{theorem}
\begin{proof}
For a maximizing $T$ and any $G\in C(X)$,
\[
 M(G)\geq\mathfrak D_G(T)
       =M(F)\int_X\mathrm{e}^{(G-F)/n}\,\mathrm{d}\nu_{T,F}.
\]
Since $B_0=-n\log M$, taking $-n$ times the logarithm and applying Jensen's inequality gives
\[
 B_0(G)\leq B_0(F)-n\log\int_X\mathrm{e}^{(G-F)/n}\,\mathrm{d}\nu_{T,F}
          \leq B_0(F)-\int_X(G-F)\,\mathrm{d}\nu_{T,F}.
\]
Hence every probability $\nu_{T,F}$ with $T\in\mathscr M(F)$ belongs to $\mathcal A_0(F)$.

Denote the set of these probabilities by $\mathscr N(F)$; by the preceding lemma it is compact and convex. To prove the reverse inclusion we compute the right derivative $(B_0)'_+(F;h)=\lim_{t\downarrow0}(B_0(F+th)-B_0(F))/t$ for $h\in C(X)$. Since $\mathfrak D_{F+th}(T)=\int_X\mathrm{e}^{th/n}\,\mathrm{d}\Gamma_F(T)$ and the masses of the measures $\Gamma_F(T)$ are bounded on $\mathscr G_\omega$, a Taylor expansion of the exponential gives, as $t\downarrow0$,
\[
 \mathfrak D_{F+th}(T)
 =\mathfrak D_F(T)+\frac tn\int_Xh\,\mathrm{d}\Gamma_F(T)+O(t^2),
\]
with the error bounded uniformly in $T$. Testing maximizers at $F$ gives
\[
 \liminf_{t\downarrow0}\frac{M(F+th)-M(F)}t
    \geq\frac{M(F)}n\max_{\nu\in\mathscr N(F)}\int_Xh\,\mathrm{d}\nu.
\]
For the reverse inequality, choose a maximizer $T_t$ at $F+th$. The multiplicative bounds in the previous proof imply $\mathfrak D_F(T_t)\to M(F)$. Every sequence $t\downarrow0$ has a subsequence along which $T_t$ converges, and the preceding lemma identifies the limit of $\Gamma_F(T_t)$ with the determinant measure of a maximizer at $F$. Since $\mathfrak D_F(T_t)\leq M(F)$, the Taylor expansion gives the matching upper bound for the difference quotients of $M$. Since $B_0=-n\log M$, therefore
\[
 (B_0)'_+(F;h)
       =-\max_{\nu\in\mathscr N(F)}\int_Xh\,\mathrm{d}\nu.
\]
If $\nu\in\mathcal A_0(F)$, its supporting inequality in direction $th$ implies
\[
 \int_Xh\,\mathrm{d}\nu\leq\max_{\widetilde\nu\in\mathscr N(F)}
                                    \int_Xh\,\mathrm{d}\widetilde\nu
 \quad\text{for every }h\in C(X).
\]
If $\nu$ were not in the weakly compact convex set $\mathscr N(F)$, the Hahn--Banach separation theorem in the dual of $C(X)$ would give some $h\in C(X)$ violating this inequality. Hence $\nu\in\mathscr N(F)$, and the two sets are equal. The absolute continuity of their members was proved in \cref{thm:weak-adjoint-probabilities}.
\end{proof}

For continuous $F$, the supporting probability is therefore unique exactly when all maximizing currents have the same normalized determinant measure, which is weaker than uniqueness of the maximizing current. Whether a maximizing current can be written as a conformal multiple of a power $(\omega+\ddc u)^{n-1}$ of a weak solution form is a further question, discussed in \cref{sec:history-gaud-volume}; the construction above works directly with the coefficient measures of the currents.
